% !PRINTED 220
\weberpage{220}

Die Gruppen $\mathfrak{A}_{\mu,\mu'}$ sind \gesperrt{konjugierte Divisoren} von $\mathfrak{A}$
(Bd.\ II, § 3), denn wenn $x_{1,0}$ durch $S$ in $x_{\mu,\mu'}$, also $R_{1,0}$ in $R_{\mu,\mu'}$
übergeht, so werden durch die Gruppe
\[
S^{-1}\mathfrak{A}_{1,0}S
\]
die Wurzeln von $R_{\mu,\mu'}$ nur unter sich vertauscht, und es ist also
\weq{6}{S^{-1}\mathfrak{A}_{1,0}S = \mathfrak{A}_{\mu,\mu'}.}

Ebenso ist, wenn $T$ eine Substitution in $\mathfrak{B}$ ist, durch die
$x_{1,0}$ in $x_{\mu,\mu'}$ übergeht:
\weq{7}{T^{-1}\mathfrak{B}_{1,0}T = \mathfrak{B}_{\mu,\mu'}.}

Der größte gemeinschaftliche Teiler $\mathfrak{A}_0$ aller Gruppen $\mathfrak{A}_{\mu,\mu'}$
wird nach (4) bestimmt durch
\[
b \equiv 0, \quad c \equiv 0, \quad a \equiv \partial \ (\mathrm{mod}\ n)
\]
und besteht also aus den Substitutionen:
\[
\begin{pmatrix} a, & 0 \\ 0, & a \end{pmatrix},
\]
worin $a$ eine beliebige, zu $n$ teilerfremde Zahl ist. $\mathfrak{A}_0$ ist identisch
mit dem größten gemeinschaftlichen Teiler dreier der Gruppen $\mathfrak{A}_{\mu,\mu'}$,
die verschiedenen Reihen angehören. Denn $\mathfrak{A}_0$ ist dann
der größte gemeinschaftliche Teiler von $\mathfrak{A}_{\mu,\mu'}$, $\mathfrak{A}_{\nu,\nu'}$, $\mathfrak{A}_{\varrho,\varrho'}$, wenn
die drei Kongruenzen
\[
\begin{aligned}
c\mu^2 - (\partial - a)\mu\mu' - b\mu'^2 &\equiv 0 \\
c\nu^2 - (\partial - a)\nu\nu' - b\nu'^2 &\equiv 0 \ (\mathrm{mod}\ n) \\
c\varrho^2 - (\partial - a)\varrho\varrho' - b\varrho'^2 &\equiv 0
\end{aligned}
\]
\gesperrt{nur} unter der Voraussetzung
\[
\partial \equiv a, \quad b \equiv 0, \quad c \equiv 0 \ (\mathrm{mod}\ n)
\]
erfüllt sind. Dies findet statt, wenn die Determinante
\[
\begin{vmatrix} \mu^2, & \mu\mu', & \mu'^2 \\ \nu^2, & \nu\nu', & \nu'^2 \\ \varrho^2, & \varrho\varrho', & \varrho'^2 \end{vmatrix}
= -(\nu\varrho' - \varrho\nu')(\varrho\mu' - \mu\varrho')(\mu\nu' - \nu\mu')
\]
relativ prim zu $n$ ist.

Ebenso ist nach (5) der größte gemeinschaftliche Teiler $\mathfrak{B}_0$
aller $\mathfrak{B}_{\mu,\mu'}$ der Inbegriff der Substitutionen:
\weq{8}{\begin{pmatrix} \alpha, & 0 \\ 0, & \alpha \end{pmatrix}}
mit der Bedingung:
\weq{9}{\alpha^2 \equiv 1 \ (\mathrm{mod}\ n).}
